\section{Estrategia OpenMP y condiciones de carrera}

\paragraph{Versión B.} Se paraleliza el ciclo externo con \code{schedule(runtime)}. Las variables del
ciclo interno se declaran dentro del ciclo, por lo que son privadas de cada hilo, y cada hilo escribe
solo la velocidad de sus propios cuerpos. El movimiento de posiciones se hace en un segundo
\code{parallel for}, después de la barrera implícita:

\begin{lstlisting}
#pragma omp parallel for schedule(runtime)
for (int i = 0; i < N; i++) {
    double fx = 0.0, fy = 0.0;            /* privadas del hilo */
    for (int j = 0; j < N; j++) { /* ... solo lee bodies[j] ... */ }
    bodies[i].vx += (fx / bodies[i].mass) * dt;  /* solo escribe su i */
    bodies[i].vy += (fy / bodies[i].mass) * dt;
}   /* barrera implícita */
#pragma omp parallel for schedule(runtime)
for (int i = 0; i < N; i++) { bodies[i].x += bodies[i].vx * dt; /* ... */ }
\end{lstlisting}

\paragraph{Versión C.} Se probaron tres formas de resolver la carrera sobre
\code{fx[i]}, \code{fy[i]}, \code{fx[j]} y \code{fy[j]}, compiladas a partir del mismo código fuente
(\code{-DSYNC\_ATOMIC}, \code{-DSYNC\_CRITICAL} o, por defecto, \code{reduction}). La
tabla~\ref{tab:sync} resume los resultados.

\begin{table}[H]
\centering
\caption{Estrategias de sincronización en la versión C (8 hilos, \code{dynamic}, chunk 64).}
\label{tab:sync}
\begin{tabular}{@{}lrr@{}}
\toprule
Estrategia & Tiempo (s) & Speedup \\
\midrule
\tabinput generado/tabla_sync.tex
\bottomrule
\end{tabular}
\end{table}

Con \code{atomic} se realizan cerca de 50 millones de operaciones atómicas por paso
($12.5\times10^6$ pares $\times$ 4), y el programa resulta más lento que el secuencial;
\code{critical} serializa por completo el núcleo del cálculo. Se \textbf{eligió \code{reduction}}
porque elimina la contención: cada hilo acumula en una copia privada de \code{fx} y \code{fy}, que se
suman al final. Su costo es memoria adicional ($2N$ \code{double} por hilo, unos 80\,KB) y una fase de
combinación $O(N\cdot p)$, ambos despreciables frente al trabajo $O(N^2)$. La diferencia numérica
máxima frente a la versión secuencial es $2.2\times10^{-8}$.
